\subsection{Вязкое трение}
%%%%%%%%%%%%%%%%%%%%%%%%%%%%%%%%%%%%%%%%%%%%%%%%%%%%
%%%%%%%%%%%%%%%%%%%%%%%%%%%%%%%%%%%%%%%%%%%%%%%%%%%%
\z{Аэростат массы $M$ опускается с постоянной скоростью. После выбрасывания груза массы $m$ аэростат стал подниматься с той же скоростью. Найдите архимедову силу, действующую на аэростат. Ускорение свободного падения равно $g$.}
%
{$F = \left(M - \frac{m}{2}\right)g$}

%%%%%%%%%%%%%%%%%%%%%%%%%%%%%%%%%%%%%%%%%%%%%%%%%%%%

 \z{Шарик массой $m$ и объёмом $V$ под действием силы тяжести падает в жидкости плотностью $\rho$ с постоянной скоростью $v$. Сила сопротивления жидкости движению шарика пропорциональна квадрату скорости. К шарику прилагается дополнительно горизонтально направленная сила $f$. Какой станет вертикальная составляющая скорости шарика $v_1$?}
%
{$v_1 = \frac{v}{\sqrt[4]{1 + \frac{f^2}{(m - \rho V)^2 g^2}}}$}
%%%%%%%%%%%%%%%%%%%%%%%%%%%%%%%%%%%%%%%%%%%%%%%%%%%%
\z{Деревянный шарик, опущенный под воду, всплывает в установившемся режиме со скоростью $v_1$, а точно такой же по размеру пластмассовый тонет со скоростью $v_2$. Куда и с какой скоростью будут двигаться в воде эти шарики, если их соединить ниткой? Сила сопротивления пропорциональна скорости, гидродинамическим взаимодействием шариков можно пренебречь. Считайте, что на движущийся шарик действует такая же сила Архимеда, как и на покоящийся.}
%
{$v = \frac{v_1 - v_2}{2}$}

%%%%%%%%%%%%%%%%%%%%%%%%%%%%%%%%%%%%%%%%%%%%%%%%%%%%	
\z{Футбольный мяч при движении в воздухе испытывает силу сопротивления, пропорциональную квадрату скорости мяча относительно воздуха. Перед ударом футболиста мяч двигался в воздухе горизонтально со скоростью $v_1 = 20\,\text{м/с}$ и с ускорением $a_1 = 13\,\text{м/с}^2$. После удара мяч полетел вертикально вверх со скоростью $v_2 = 10\,\text{м/с}$. Каково ускорение мяча сразу после удара?}
%
{$a_2 = g + \left(\frac{v_2}{v_1}\right)^2 \sqrt{a_1^2 - g^2} \approx 11.9 \, \text{м/с}^2 $}

%%%%%%%%%%%%%%%%%%%%%%%%%%%%%%%%%%%%%%%%%%%%%%%%%%%%
\z{Лодку оттолкнули от берега озера, сообщив ей скорость $v_0 = 1\,\text{м/с}$. Лодка, двигаясь прямолинейно, имела на расстоянии $s_1 = 14\,\text{м}$ от берега скорость $v_1 = 0.3\,\text{м/с}$. На каком расстоянии от берега скорость лодки была $v = 0.5\,\text{м/с}$? Считать, что сила сопротивления движению лодки пропорциональна её скорости.}
%
{$s = s_1 \frac{v - v_0}{v_1 - v_0} = 10 \, \text{м} $}

